% =============================================================================
% Project : THE PHILOSOPHY OF VIABILITY
% Book    : Book I — Ontology of Viability
% Book ID : VS.B01
% Part    : Part I — Identity and Manifestation
% Part ID : VS.B01.P01
% Chapter : Chapter 08 — Comparative Philosophical Examination
% Chapter ID : VS.B01.P01.C08
%
% File :
%   4_ASSETS/diagrams/
%   VS.B01.P01.C08.S12.SS09.P004.DG01.al_farabi_emanative_cosmological_order.tex
%
% Section : Section 12 — Al-Farabi
%
% Asset :
%   Diagram DG01 — Al-Farabi's Emanative Cosmological Order
%
% Structural Identifier :
%   VS.B01.P01.C08.S12.SS09.P004.DG01
%
% Purpose :
%   To represent the broad causal architecture of Al-Farabi's
%   emanative cosmology and distinguish the First Existent,
%   the secondary intellects, the celestial order, the Active
%   Intellect, and the sublunary domain.
%
% Intellectual Benefit :
%   Clarifies the broad relations of causal dependence while
%   avoiding an unsupported one-to-one mapping of every
%   intellect to a celestial sphere across all relevant texts.
%
% Aesthetic / Experiential Benefit :
%   Provides a compact visual synthesis of a complex cosmological
%   structure.
%
% Representational Limits :
%   1. The diagram is schematic, not exhaustive.
%   2. It does not represent a temporal sequence.
%   3. The arrangement of immediate effects differs between
%      the Perfect State and the Political Regime.
%   4. The Active Intellect is shown separately to clarify
%      its function; it belongs to the relevant series of
%      separate intellects.
%   5. The diagram is not a representation of the domains
%      or hubs proposed in Book I.
%   6. The Active Intellect is not identified with Light.
%
% Dependencies :
%   TikZ; positioning library; arrows.meta. (no other libraries required;
%   all widths scale with \linewidth).
%
% =============================================================================

\begin{figure}[htbp]
\centering

\begin{tikzpicture}[
    font=\small,
    node distance=0.95cm,
    box/.style={
        rounded corners=4pt,
        align=center,
        inner sep=6pt,
        minimum height=1.2cm
    },
    first/.style={
        box,
        draw=orange!75!black,
        fill=yellow!15,
        line width=1.0pt,
        text width=0.50\linewidth
    },
    intellect/.style={
        box,
        draw=blue!75!black,
        fill=blue!8,
        line width=0.8pt,
        text width=0.50\linewidth
    },
    active/.style={
        box,
        draw=green!55!black,
        fill=green!9,
        line width=0.9pt,
        text width=0.50\linewidth
    },
    celestial/.style={
        box,
        draw=blue!65!black,
        fill=gray!8,
        line width=0.7pt,
        text width=0.26\linewidth
    },
    sublunary/.style={
        box,
        draw=blue!65!black,
        fill=gray!15,
        line width=0.8pt,
        text width=0.50\linewidth
    },
    note/.style={
        draw=gray!65,
        fill=gray!5,
        rounded corners=3pt,
        text width=0.80\linewidth,
        align=left,
        font=\footnotesize,
        inner sep=6pt
    },
    relation/.style={
        draw=blue!65!black,
        line width=0.85pt,
        -{Stealth[length=2mm]}
    },
    explanatory/.style={
        draw=gray!65,
        line width=0.7pt,
        dashed,
        -{Stealth[length=1.8mm]}
    },
    elabel/.style={
        font=\scriptsize,
        align=left,
        inner sep=2pt
    }
]

% ---------------------------------------------------------------------------
% Main vertical axis: First Existent -> Secondary Intellects -> Active Intellect
% ---------------------------------------------------------------------------

\node[first] (first) {
    \textbf{The First Existent / First Cause}\\
    Uncaused, independent, and simple
};

\node[intellect, below=of first] (secondary) {
    \textbf{Secondary Separate Intellects}\\
    An ordered series of intellectual principles
};

\node[active, below=of secondary] (active) {
    \textbf{Active Intellect}\\
    Terminal separate intellect in the familiar sublunary scheme
};

\node[sublunary, below=1.15cm of active] (sublunary) {
    \textbf{Sublunary Domain}\\
    Material beings, forms, souls, generation, and corruption
};

% ---------------------------------------------------------------------------
% Celestial Order (beside the series, outside it)
% ---------------------------------------------------------------------------

\node[celestial, right=1.7cm of secondary] (spheres) {
    \textbf{Celestial Order}\\
    Spheres and their associated order
};

% ---------------------------------------------------------------------------
% Frame: the Active Intellect belongs to the series of separate intellects
% ---------------------------------------------------------------------------

\draw[blue!60!black, densely dotted, line width=0.6pt, rounded corners=6pt]
    ([xshift=-3.5mm, yshift=9mm]secondary.north west) rectangle
    ([xshift= 3.5mm, yshift=-3.5mm]active.south east);

\node[font=\scriptsize\itshape, text=blue!60!black, anchor=north west,
      inner sep=2pt]
    at ([xshift=-3.5mm, yshift=9mm]secondary.north west)
    {Series of separate intellects};

% ---------------------------------------------------------------------------
% Relations
% ---------------------------------------------------------------------------

\draw[relation] (first.south) -- 
    node[elabel, right, xshift=2pt] {Causal dependence}
    (secondary.north);

\draw[relation] (secondary.south) -- 
    node[elabel, right, xshift=2pt] {Continued ordered dependence}
    (active.north);

\draw[explanatory] (secondary.east) -- 
    node[elabel, above, align=center] {associated\\with}
    (spheres.west);

\draw[relation] (active.south) -- 
    node[elabel, right, xshift=2pt] {Governing and formative relation}
    (sublunary.north);

% ---------------------------------------------------------------------------
% Interpretive limits
% ---------------------------------------------------------------------------

\node[note, anchor=north west]
    (limits) at ([xshift=-3.5mm, yshift=-7mm]sublunary.south west) {
    \textbf{Interpretive limits.} Schematic only: no uniform one-to-one
    correspondence between each separate intellect and each celestial
    sphere is claimed across all treatises. The Active Intellect is
    displayed apart for clarity, not as a principle outside the intellectual
    series. Relations are metaphysical and causal, not temporal stages.
};

\end{tikzpicture}

\caption{Schematic representation of the broad causal architecture of
Al-Farabi's emanative cosmology. The First Existent is the first cause of
the ordered series of existents. The secondary intellects are associated
with the celestial order, while the Active Intellect, as the terminal member
of the series, stands in a governing and formative relation to the sublunary
domain. The diagram does not imply that the relevant treatises present an
identical arrangement.}

\label{fig:VS.B01.P01.C08.S12.SS09.P004.DG01}

\end{figure}

% =============================================================================
% END OF DIAGRAM DG01
% =============================================================================